\documentclass[11pt,a4paper]{article}
\usepackage{curso}
\begin{document}
\portada{Lección}{8.5}{K-FAC termodinámico}

\noindent\textbf{Unidad 8.}\quad\textbf{Objetivos:} reemplazar las inversiones $n\times n$ de K-FAC por soluciones termodinámicas; comparar complejidades;
cuantización de entrada y salida.
\textbf{Prerrequisitos:} 5.6, 6.1--6.2, 8.4.

\section{El algoritmo}
La GPU calcula $A=\mathbb E[aa^\top]$ y $G=\mathbb E[gg^\top]$ por capa; en vez de invertirlas digitalmente, las carga (DAC) en un resolvedor termodinámico de
sistemas de tamaño $n$ (dinámica OU, Unidad 6), lee la solución (ADC) y aplica $\Delta W=-\eta\,G^{-1}\nabla W A^{-1}$. Los promedios móviles pueden vivir en la GPU o en
el dispositivo.

\section{Complejidad (Donatella et al.\ 2025, por capa y por iteración)}
\begin{center}\begin{tabular}{lll}\toprule
optimizador & tiempo & memoria\\\midrule
SGD / Adam & $O(bn^2)$ & $O(n^2)$\\
K-FAC & $O(bn^2+n^3)$ & $O(bn+n^2)$\\
K-FAC termodinámico & $O(bn^2+n^2\kappa^2)$ & $O(bn)$ (sin EMA)\\\bottomrule\end{tabular}\end{center}
$\kappa$ acota el número de condición de los factores. La ventaja crece linealmente con el ancho $n$.

\section{Aceleraciones estimadas}
Fracción del tiempo de K-FAC dedicada a invertir (AlgoPerf, medida en GPU): 11\,\% (ViT/ImageNet), 27\,\% (GNN/ogbg-molpcba). Los autores estiman alcanzar la misma métrica de
validación $\approx50\,\%$ más rápido en la GNN \evP. Por Amdahl, acelerar solo la inversión con el mismo número de pasos da a lo sumo $1/(1-0{,}27)=1{,}37$: la estimación debe
incluir cambios en los pasos. Ningún dispositivo físico ejecutó estos experimentos.

\section{Cuantización}
Entrada (DAC): un cuantizador ingenuo puede romper la semidefinición positiva; se usa uno conservador de diagonal dominante. Salida (ADC): en el entrenamiento de una
ResNet, la entrada afecta poco y la salida necesita al menos 8 bits; con 8 bits, K-FAC compite con Adam en precisión completa \evS.

\begin{ejemplo}[ (advertencia metodológica)]
En la capa softmax de 8.4, un solo paso \evS: con cuantizador uniforme ingenuo, la entrada daña más (coseno $0{,}70$ a 4 bits) y la salida apenas cambia la dirección
($0{,}93$). Es lo contrario del hallazgo en entrenamiento completo: las métricas de un paso no sustituyen experimentos de entrenamiento.
\end{ejemplo}

\begin{aplicacion}
Redes anchas; medir primero la fracción de inversión; cuantizador conservador y ADC de $\ge8$ bits. Laboratorio L11: barrido de bits de entrada y salida.
\end{aplicacion}

\begin{error}
\begin{itemize}[nosep]
\item Prometer aceleración total cuando la inversión es una fracción. \quad Ignorar $\kappa^2$.
\item Cuantizadores que rompen la positividad. \quad Juzgar la precisión con métricas de un paso.
\end{itemize}
\end{error}

\begin{thebibliography}{9}\small
\bibitem[41]{stso} K. Donatella et al., Scalable thermodynamic second-order optimization, arXiv:2502.08603 (2025).
\bibitem{mg} J. Martens y R. Grosse, Optimizing neural networks with Kronecker-factored approximate curvature, \emph{ICML} (2015).
\bibitem[30]{tla} M. Aifer et al., Thermodynamic linear algebra, \emph{npj Unconv. Comput.} 1, 13 (2024).
\end{thebibliography}
\end{document}
